\section{十九世纪冲突证据链事件锚点}
以下锚点追踪内外战争的文件、条件、决策与后果，提示战场、城市、省级财政和家庭生活的材料并不来自同一种档案。
\subsection{1850—1859}
\eventanchor{EQ-1851-JINTIAN-UPRISING-DOCUMENTS}{清廷围绕「洪秀全集团在金田起事…}咸丰元年（1851年），清廷围绕「洪秀全集团在金田起事，太平天国战争开始」调遣军队、筹措饷械并处理战报，中央与地方的军政文书形成连续记录。核心取证：《清文宗实录》咸丰元年相关条；军机处档、战事方略及地方奏报。这一锚点记录的是文书链：它能证明呈报、上谕、部议或照会进入机构流程，不能由此推出实际执行。来源保留：此行仅确认相关文书链和可追查的行政动作；不得由其直接推断政策有效、地方服从、民众态度或单一因果。
\eventanchor{EQ-1851-YONGAN-KINGSHIPS-DOCUMENTS}{围绕「太平军在永安建立王号和军…}咸丰元年（1851年），围绕「太平军在永安建立王号和军事组织」的上谕、奏折与部议在中央—地方行政链中流转。核心取证：《清文宗实录》咸丰元年相关条；宫中朱批奏折、内阁大库题本与《清史稿》相关志传。这一锚点记录的是文书链：它能证明呈报、上谕、部议或照会进入机构流程，不能由此推出实际执行。来源保留：此行仅确认相关文书链和可追查的行政动作；不得由其直接推断政策有效、地方服从、民众态度或单一因果。
\eventanchor{EQ-1852-TAIPING-NORTHWARD-MARCH-DOCUMENTS}{清廷围绕「太平军北上经湖南，清…}咸丰二年（1852年），清廷围绕「太平军北上经湖南，清军防御与地方团练开始加速组织」调遣军队、筹措饷械并处理战报，中央与地方的军政文书形成连续记录。核心取证：《清文宗实录》咸丰二年相关条；军机处档、战事方略及地方奏报。这一锚点记录的是文书链：它能证明呈报、上谕、部议或照会进入机构流程，不能由此推出实际执行。来源保留：此行仅确认相关文书链和可追查的行政动作；不得由其直接推断政策有效、地方服从、民众态度或单一因果。
\eventanchor{EQ-1852-HUNAN-ARMY-ORIGINS-DOCUMENTS}{清廷围绕「曾国藩等在湖南组织团…}咸丰二年（1852年），清廷围绕「曾国藩等在湖南组织团练和湘军雏形」调遣军队、筹措饷械并处理战报，中央与地方的军政文书形成连续记录。核心取证：《清文宗实录》咸丰二年相关条；军机处档、战事方略及地方奏报。这一锚点记录的是文书链：它能证明呈报、上谕、部议或照会进入机构流程，不能由此推出实际执行。来源保留：此行仅确认相关文书链和可追查的行政动作；不得由其直接推断政策有效、地方服从、民众态度或单一因果。
\eventanchor{EQ-1853-NANJING-TAKING-DOCUMENTS}{清廷围绕「太平军攻占南京并改名…}咸丰三年（1853年），清廷围绕「太平军攻占南京并改名天京，建立长期政权中心」调遣军队、筹措饷械并处理战报，中央与地方的军政文书形成连续记录。核心取证：《清文宗实录》咸丰三年相关条；军机处档、战事方略及地方奏报。这一锚点记录的是文书链：它能证明呈报、上谕、部议或照会进入机构流程，不能由此推出实际执行。来源保留：此行仅确认相关文书链和可追查的行政动作；不得由其直接推断政策有效、地方服从、民众态度或单一因果。
\eventanchor{EQ-1853-NORTHERN-EXPEDITION-DOCUMENTS}{清廷围绕「太平军发动北伐，进逼…}咸丰三年（1853年），清廷围绕「太平军发动北伐，进逼华北但补给与指挥困难加重」调遣军队、筹措饷械并处理战报，中央与地方的军政文书形成连续记录。核心取证：《清文宗实录》咸丰三年相关条；军机处档、战事方略及地方奏报。这一锚点记录的是文书链：它能证明呈报、上谕、部议或照会进入机构流程，不能由此推出实际执行。来源保留：此行仅确认相关文书链和可追查的行政动作；不得由其直接推断政策有效、地方服从、民众态度或单一因果。
\eventanchor{EQ-1853-WESTERN-EXPEDITION-DOCUMENTS}{清廷围绕「太平军西征争夺长江中…}咸丰三年（1853年），清廷围绕「太平军西征争夺长江中上游，战事牵动安徽、江西、湖北」调遣军队、筹措饷械并处理战报，中央与地方的军政文书形成连续记录。核心取证：《清文宗实录》咸丰三年相关条；军机处档、战事方略及地方奏报。这一锚点记录的是文书链：它能证明呈报、上谕、部议或照会进入机构流程，不能由此推出实际执行。来源保留：此行仅确认相关文书链和可追查的行政动作；不得由其直接推断政策有效、地方服从、民众态度或单一因果。
\eventanchor{EQ-1853-SMALL-SWORDS-SHANGHAI-DOCUMENTS}{围绕「上海小刀会起事，租界和华…}咸丰三年（1853年），围绕「上海小刀会起事，租界和华人城市空间关系发生变化」的地方呈报、案件或赈务记录将社会反应带入官府文书链。核心取证：《清文宗实录》咸丰三年相关条；地方志、督抚奏折及地方档案。这一锚点记录的是文书链：它能证明呈报、上谕、部议或照会进入机构流程，不能由此推出实际执行。来源保留：此行仅确认相关文书链和可追查的行政动作；不得由其直接推断政策有效、地方服从、民众态度或单一因果。
\eventanchor{EQ-1854-TAIPING-NORTHERN-DEFEAT-DOCUMENTS}{清廷围绕「太平军北伐受挫，清军…}咸丰四年（1854年），清廷围绕「太平军北伐受挫，清军逐步恢复华北防线」调遣军队、筹措饷械并处理战报，中央与地方的军政文书形成连续记录。核心取证：《清文宗实录》咸丰四年相关条；军机处档、战事方略及地方奏报。这一锚点记录的是文书链：它能证明呈报、上谕、部议或照会进入机构流程，不能由此推出实际执行。来源保留：此行仅确认相关文书链和可追查的行政动作；不得由其直接推断政策有效、地方服从、民众态度或单一因果。
\eventanchor{EQ-1854-XIANG-ARMY-ORGANIZATION-DOCUMENTS}{围绕「湘军编制、饷源和将领网络…}咸丰四年（1854年），围绕「湘军编制、饷源和将领网络逐步成形」的上谕、奏折与部议在中央—地方行政链中流转。核心取证：《清文宗实录》咸丰四年相关条；宫中朱批奏折、内阁大库题本与《清史稿》相关志传。这一锚点记录的是文书链：它能证明呈报、上谕、部议或照会进入机构流程，不能由此推出实际执行。来源保留：此行仅确认相关文书链和可追查的行政动作；不得由其直接推断政策有效、地方服从、民众态度或单一因果。
\eventanchor{EQ-1855-YELLOW-RIVER-COURSE-SHIFT-DOCUMENTS}{围绕「黄河在铜瓦厢决口改道北流…}咸丰五年（1855年），围绕「黄河在铜瓦厢决口改道北流，重塑华北水系与社会经济」的地方呈报、案件或赈务记录将社会反应带入官府文书链。核心取证：《清文宗实录》咸丰五年相关条；地方志、督抚奏折及地方档案。这一锚点记录的是文书链：它能证明呈报、上谕、部议或照会进入机构流程，不能由此推出实际执行。来源保留：此行仅确认相关文书链和可追查的行政动作；不得由其直接推断政策有效、地方服从、民众态度或单一因果。
\eventanchor{EQ-1855-NIAN-REBELLION-EXPANSION-DOCUMENTS}{清廷围绕「捻军在华北、安徽等地…}咸丰五年（1855年），清廷围绕「捻军在华北、安徽等地活动扩大，形成与太平战争交织的战场」调遣军队、筹措饷械并处理战报，中央与地方的军政文书形成连续记录。核心取证：《清文宗实录》咸丰五年相关条；军机处档、战事方略及地方奏报。这一锚点记录的是文书链：它能证明呈报、上谕、部议或照会进入机构流程，不能由此推出实际执行。来源保留：此行仅确认相关文书链和可追查的行政动作；不得由其直接推断政策有效、地方服从、民众态度或单一因果。
\eventanchor{EQ-1856-TIANJING-INCIDENT-DOCUMENTS}{围绕「天京事变中太平天国领导层…}咸丰六年（1856年），围绕「天京事变中太平天国领导层发生内讧和大规模杀戮」的上谕、奏折与部议在中央—地方行政链中流转。核心取证：《清文宗实录》咸丰六年相关条；宫中朱批奏折、内阁大库题本与《清史稿》相关志传。这一锚点记录的是文书链：它能证明呈报、上谕、部议或照会进入机构流程，不能由此推出实际执行。来源保留：此行仅确认相关文书链和可追查的行政动作；不得由其直接推断政策有效、地方服从、民众态度或单一因果。
\eventanchor{EQ-1856-SECOND-OPIUM-WAR-BEGIN-DOCUMENTS}{清廷围绕「亚罗号事件等冲突后英…}咸丰六年（1856年），清廷围绕「亚罗号事件等冲突后英法对清开战，第二次鸦片战争开始」调遣军队、筹措饷械并处理战报，中央与地方的军政文书形成连续记录。核心取证：《清文宗实录》咸丰六年相关条；军机处档、战事方略及地方奏报。这一锚点记录的是文书链：它能证明呈报、上谕、部议或照会进入机构流程，不能由此推出实际执行。来源保留：此行仅确认相关文书链和可追查的行政动作；不得由其直接推断政策有效、地方服从、民众态度或单一因果。
\eventanchor{EQ-1857-GUANGZHOU-FALL-DOCUMENTS}{清廷围绕「英法联军攻占广州，叶…}咸丰七年（1857年），清廷围绕「英法联军攻占广州，叶名琛被俘」调遣军队、筹措饷械并处理战报，中央与地方的军政文书形成连续记录。核心取证：《清文宗实录》咸丰七年相关条；军机处档、战事方略及地方奏报。这一锚点记录的是文书链：它能证明呈报、上谕、部议或照会进入机构流程，不能由此推出实际执行。来源保留：此行仅确认相关文书链和可追查的行政动作；不得由其直接推断政策有效、地方服从、民众态度或单一因果。
\eventanchor{EQ-1858-TIANJIN-TREATIES-DOCUMENTS}{围绕「《天津条约》等在战争压力…}咸丰八年（1858年），围绕「《天津条约》等在战争压力下签订，扩大通商、使节与传教权利」的照会、条约文本、换约或地方执行报告分别进入外交文书链。核心取证：《清文宗实录》咸丰八年相关条；《筹办夷务始末》相关年卷与中外照会、条约文本。这一锚点记录的是文书链：它能证明呈报、上谕、部议或照会进入机构流程，不能由此推出实际执行。来源保留：此行仅确认相关文书链和可追查的行政动作；不得由其直接推断政策有效、地方服从、民众态度或单一因果。
\eventanchor{EQ-1859-TAKU-BATTLE-DOCUMENTS}{清廷围绕「大沽口战事使英法换约…}咸丰九年（1859年），清廷围绕「大沽口战事使英法换约受阻，战争再度升级」调遣军队、筹措饷械并处理战报，中央与地方的军政文书形成连续记录。核心取证：《清文宗实录》咸丰九年相关条；军机处档、战事方略及地方奏报。这一锚点记录的是文书链：它能证明呈报、上谕、部议或照会进入机构流程，不能由此推出实际执行。来源保留：此行仅确认相关文书链和可追查的行政动作；不得由其直接推断政策有效、地方服从、民众态度或单一因果。
\eventanchor{EQ-1851-JINTIAN-UPRISING-PRELUDE}{「洪秀全集团在金田起事，太平天…}咸丰元年（1851年），「洪秀全集团在金田起事，太平天国战争开始」之前的条件、信息和争议已通过中央与地方材料显现，构成该事件可追踪的前史节点。核心取证：《清文宗实录》咸丰元年相关条；军机处档、战事方略及地方奏报；同时期地方档案、地方志、日记或对方材料应按具体地区补检。这一锚点界定可回查的前史条件；条件、传闻或信息累积不应被写成已经证实的单一直接原因。来源保留：官修编年和地方材料的记录密度与立场并不相同；本行不将条件、传闻、呈报或后见解释直接等同为确定因果。
\eventanchor{EQ-1851-JINTIAN-UPRISING-DECISION}{围绕「洪秀全集团在金田起事，太…}咸丰元年（1851年），围绕「洪秀全集团在金田起事，太平天国战争开始」的命令、调度、照会或呈报形成独立的中央—地方行动文书链。核心取证：《清文宗实录》咸丰元年相关条；军机处档、战事方略及地方奏报。这一锚点定位决策或调度动作；命令发出、地方接收、执行过程和实际结果仍须分别寻找证据。来源保留：奏折、上谕、部议、军机档或条约可证明行政动作和机构立场，不自动证明所有地区、群体或对手按其意图行动。
\eventanchor{EQ-1851-JINTIAN-UPRISING-AFTERMATH}{「洪秀全集团在金田起事，太平天…}咸丰元年（1851年），「洪秀全集团在金田起事，太平天国战争开始」引发的善后、执行或地方回应构成可由不同材料追踪的后续节点。核心取证：《清文宗实录》咸丰元年相关条；军机处档、战事方略及地方奏报；相关地区的档案、志书、契约、报刊或对方材料。这一锚点追踪后续影响；不同地区、群体与时段的反应不能由战后或改革后的概括性记忆代替。来源保留：战后、改革后或条约后的记忆、地方记录和官方总结常具有事后选择性；后续影响须按地点、群体和时间分层，不作整体化推断。
\eventanchor{EQ-1851-YONGAN-KINGSHIPS-PRELUDE}{「太平军在永安建立王号和军事组…}咸丰元年（1851年），「太平军在永安建立王号和军事组织」之前的条件、信息和争议已通过中央与地方材料显现，构成该事件可追踪的前史节点。核心取证：《清文宗实录》咸丰元年相关条；宫中朱批奏折、内阁大库题本与《清史稿》相关志传；同时期地方档案、地方志、日记或对方材料应按具体地区补检。这一锚点界定可回查的前史条件；条件、传闻或信息累积不应被写成已经证实的单一直接原因。来源保留：官修编年和地方材料的记录密度与立场并不相同；本行不将条件、传闻、呈报或后见解释直接等同为确定因果。
\eventanchor{EQ-1851-YONGAN-KINGSHIPS-DECISION}{围绕「太平军在永安建立王号和军…}咸丰元年（1851年），围绕「太平军在永安建立王号和军事组织」的命令、调度、照会或呈报形成独立的中央—地方行动文书链。核心取证：《清文宗实录》咸丰元年相关条；宫中朱批奏折、内阁大库题本与《清史稿》相关志传。这一锚点定位决策或调度动作；命令发出、地方接收、执行过程和实际结果仍须分别寻找证据。来源保留：奏折、上谕、部议、军机档或条约可证明行政动作和机构立场，不自动证明所有地区、群体或对手按其意图行动。
\eventanchor{EQ-1851-YONGAN-KINGSHIPS-AFTERMATH}{「太平军在永安建立王号和军事组…}咸丰元年（1851年），「太平军在永安建立王号和军事组织」引发的善后、执行或地方回应构成可由不同材料追踪的后续节点。核心取证：《清文宗实录》咸丰元年相关条；宫中朱批奏折、内阁大库题本与《清史稿》相关志传；相关地区的档案、志书、契约、报刊或对方材料。这一锚点追踪后续影响；不同地区、群体与时段的反应不能由战后或改革后的概括性记忆代替。来源保留：战后、改革后或条约后的记忆、地方记录和官方总结常具有事后选择性；后续影响须按地点、群体和时间分层，不作整体化推断。
\eventanchor{EQ-1852-HUNAN-ARMY-ORIGINS-PRELUDE}{「曾国藩等在湖南组织团练和湘军…}咸丰二年（1852年），「曾国藩等在湖南组织团练和湘军雏形」之前的条件、信息和争议已通过中央与地方材料显现，构成该事件可追踪的前史节点。核心取证：《清文宗实录》咸丰二年相关条；军机处档、战事方略及地方奏报；同时期地方档案、地方志、日记或对方材料应按具体地区补检。这一锚点界定可回查的前史条件；条件、传闻或信息累积不应被写成已经证实的单一直接原因。来源保留：官修编年和地方材料的记录密度与立场并不相同；本行不将条件、传闻、呈报或后见解释直接等同为确定因果。
\eventanchor{EQ-1852-HUNAN-ARMY-ORIGINS-DECISION}{围绕「曾国藩等在湖南组织团练和…}咸丰二年（1852年），围绕「曾国藩等在湖南组织团练和湘军雏形」的命令、调度、照会或呈报形成独立的中央—地方行动文书链。核心取证：《清文宗实录》咸丰二年相关条；军机处档、战事方略及地方奏报。这一锚点定位决策或调度动作；命令发出、地方接收、执行过程和实际结果仍须分别寻找证据。来源保留：奏折、上谕、部议、军机档或条约可证明行政动作和机构立场，不自动证明所有地区、群体或对手按其意图行动。
\eventanchor{EQ-1852-HUNAN-ARMY-ORIGINS-AFTERMATH}{「曾国藩等在湖南组织团练和湘军…}咸丰二年（1852年），「曾国藩等在湖南组织团练和湘军雏形」引发的善后、执行或地方回应构成可由不同材料追踪的后续节点。核心取证：《清文宗实录》咸丰二年相关条；军机处档、战事方略及地方奏报；相关地区的档案、志书、契约、报刊或对方材料。这一锚点追踪后续影响；不同地区、群体与时段的反应不能由战后或改革后的概括性记忆代替。来源保留：战后、改革后或条约后的记忆、地方记录和官方总结常具有事后选择性；后续影响须按地点、群体和时间分层，不作整体化推断。
\eventanchor{EQ-1853-NANJING-TAKING-PRELUDE}{「太平军攻占南京并改名天京，建…}咸丰三年（1853年），「太平军攻占南京并改名天京，建立长期政权中心」之前的条件、信息和争议已通过中央与地方材料显现，构成该事件可追踪的前史节点。核心取证：《清文宗实录》咸丰三年相关条；军机处档、战事方略及地方奏报；同时期地方档案、地方志、日记或对方材料应按具体地区补检。这一锚点界定可回查的前史条件；条件、传闻或信息累积不应被写成已经证实的单一直接原因。来源保留：官修编年和地方材料的记录密度与立场并不相同；本行不将条件、传闻、呈报或后见解释直接等同为确定因果。
\subsection{1860—1869}
\eventanchor{EQ-1860-BEIJING-CAMPAIGN-DOCUMENTS}{清廷围绕「英法联军进逼北京，圆…}咸丰十年（1860年），清廷围绕「英法联军进逼北京，圆明园遭焚掠，咸丰帝出走热河」调遣军队、筹措饷械并处理战报，中央与地方的军政文书形成连续记录。核心取证：《清文宗实录》咸丰十年相关条；军机处档、战事方略及地方奏报。这一锚点记录的是文书链：它能证明呈报、上谕、部议或照会进入机构流程，不能由此推出实际执行。来源保留：此行仅确认相关文书链和可追查的行政动作；不得由其直接推断政策有效、地方服从、民众态度或单一因果。
\eventanchor{EQ-1860-CONVENTION-OF-PEKING-DOCUMENTS}{围绕「《北京条约》及相关中俄条…}咸丰十年（1860年），围绕「《北京条约》及相关中俄条约确认和扩大列强权益并涉及东北边界」的照会、条约文本、换约或地方执行报告分别进入外交文书链。核心取证：《清文宗实录》咸丰十年相关条；《筹办夷务始末》相关年卷与中外照会、条约文本。这一锚点记录的是文书链：它能证明呈报、上谕、部议或照会进入机构流程，不能由此推出实际执行。来源保留：此行仅确认相关文书链和可追查的行政动作；不得由其直接推断政策有效、地方服从、民众态度或单一因果。
\eventanchor{EQ-1861-XIANFENG-DEATH-DOCUMENTS}{围绕「咸丰帝去世，同治帝即位，…}咸丰十一年（1861年），围绕「咸丰帝去世，同治帝即位，顾命大臣与两宫太后权力冲突爆发」的上谕、奏折与部议在中央—地方行政链中流转。核心取证：《清文宗实录》咸丰十一年相关条；宫中朱批奏折、内阁大库题本与《清史稿》相关志传。这一锚点记录的是文书链：它能证明呈报、上谕、部议或照会进入机构流程，不能由此推出实际执行。来源保留：此行仅确认相关文书链和可追查的行政动作；不得由其直接推断政策有效、地方服从、民众态度或单一因果。
\eventanchor{EQ-1861-XINYOU-COUP-DOCUMENTS}{围绕「慈禧、慈安与恭亲王发动辛…}咸丰十一年（1861年），围绕「慈禧、慈安与恭亲王发动辛酉政变，肃顺等顾命大臣失势」的上谕、奏折与部议在中央—地方行政链中流转。核心取证：《清文宗实录》咸丰十一年相关条；宫中朱批奏折、内阁大库题本与《清史稿》相关志传。这一锚点记录的是文书链：它能证明呈报、上谕、部议或照会进入机构流程，不能由此推出实际执行。来源保留：此行仅确认相关文书链和可追查的行政动作；不得由其直接推断政策有效、地方服从、民众态度或单一因果。
\eventanchor{EQ-1861-ZONGLI-YAMEN-DOCUMENTS}{围绕「总理各国事务衙门设立，清…}咸丰十一年（1861年），围绕「总理各国事务衙门设立，清廷开始形成常设近代外交处理机构」的照会、条约文本、换约或地方执行报告分别进入外交文书链。核心取证：《清文宗实录》咸丰十一年相关条；《筹办夷务始末》相关年卷与中外照会、条约文本。这一锚点记录的是文书链：它能证明呈报、上谕、部议或照会进入机构流程，不能由此推出实际执行。来源保留：此行仅确认相关文书链和可追查的行政动作；不得由其直接推断政策有效、地方服从、民众态度或单一因果。
\eventanchor{EQ-1862-SELF-STRENGTHENING-ORIGINS-DOCUMENTS}{围绕「洋务实践在军械、翻译、外…}同治元年（1862年），围绕「洋务实践在军械、翻译、外交和地方军政中展开」的上谕、奏折与部议在中央—地方行政链中流转。核心取证：《清穆宗实录》同治元年相关条；宫中朱批奏折、内阁大库题本与《清史稿》相关志传。这一锚点记录的是文书链：它能证明呈报、上谕、部议或照会进入机构流程，不能由此推出实际执行。来源保留：此行仅确认相关文书链和可追查的行政动作；不得由其直接推断政策有效、地方服从、民众态度或单一因果。
\eventanchor{EQ-1862-SHANGHAI-FOREIGN-ARMS-DOCUMENTS}{清廷围绕「上海的军火、雇佣力量…}同治元年（1862年），清廷围绕「上海的军火、雇佣力量和租界网络影响清军对太平军作战」调遣军队、筹措饷械并处理战报，中央与地方的军政文书形成连续记录。核心取证：《清穆宗实录》同治元年相关条；军机处档、战事方略及地方奏报。这一锚点记录的是文书链：它能证明呈报、上谕、部议或照会进入机构流程，不能由此推出实际执行。来源保留：此行仅确认相关文书链和可追查的行政动作；不得由其直接推断政策有效、地方服从、民众态度或单一因果。
\eventanchor{EQ-1863-EVER-VICTORIOUS-ARMY-DOCUMENTS}{清廷围绕「常胜军在江南战场参与…}同治二年（1863年），清廷围绕「常胜军在江南战场参与清军作战，戈登接掌后扩大影响」调遣军队、筹措饷械并处理战报，中央与地方的军政文书形成连续记录。核心取证：《清穆宗实录》同治二年相关条；军机处档、战事方略及地方奏报。这一锚点记录的是文书链：它能证明呈报、上谕、部议或照会进入机构流程，不能由此推出实际执行。来源保留：此行仅确认相关文书链和可追查的行政动作；不得由其直接推断政策有效、地方服从、民众态度或单一因果。
\eventanchor{EQ-1863-NIAN-WAR-MOBILIZATION-DOCUMENTS}{清廷围绕「清廷调集湘淮军和地方…}同治二年（1863年），清廷围绕「清廷调集湘淮军和地方力量应对捻军，华北战事加剧」调遣军队、筹措饷械并处理战报，中央与地方的军政文书形成连续记录。核心取证：《清穆宗实录》同治二年相关条；军机处档、战事方略及地方奏报。这一锚点记录的是文书链：它能证明呈报、上谕、部议或照会进入机构流程，不能由此推出实际执行。来源保留：此行仅确认相关文书链和可追查的行政动作；不得由其直接推断政策有效、地方服从、民众态度或单一因果。
\eventanchor{EQ-1864-TIANJING-FALL-DOCUMENTS}{清廷围绕「湘军攻陷天京，太平天…}同治三年（1864年），清廷围绕「湘军攻陷天京，太平天国核心政权覆亡」调遣军队、筹措饷械并处理战报，中央与地方的军政文书形成连续记录。核心取证：《清穆宗实录》同治三年相关条；军机处档、战事方略及地方奏报。这一锚点记录的是文书链：它能证明呈报、上谕、部议或照会进入机构流程，不能由此推出实际执行。来源保留：此行仅确认相关文书链和可追查的行政动作；不得由其直接推断政策有效、地方服从、民众态度或单一因果。
\eventanchor{EQ-1864-TAIPING-AFTERMATH}{太平战争结束后江南开始处理难民…}同治三年（1864年），太平战争结束后江南开始处理难民、田地、税收与城镇重建。核心取证：《清穆宗实录》同治三年相关条；地方志、督抚奏折及地方档案。这一锚点追踪后续影响；不同地区、群体与时段的反应不能由战后或改革后的概括性记忆代替。来源保留：以同治三年（1864年）的同期文书为定位起点；官修记录、奏报或外交文本服务特定机构立场，具体日期、规模、地方执行与对手视角须以独立材料复核。
\eventanchor{EQ-1864-TAIPING-AFTERMATH-DOCUMENTS}{围绕「太平战争结束后江南开始处…}同治三年（1864年），围绕「太平战争结束后江南开始处理难民、田地、税收与城镇重建」的地方呈报、案件或赈务记录将社会反应带入官府文书链。核心取证：《清穆宗实录》同治三年相关条；地方志、督抚奏折及地方档案。这一锚点记录的是文书链：它能证明呈报、上谕、部议或照会进入机构流程，不能由此推出实际执行。来源保留：此行仅确认相关文书链和可追查的行政动作；不得由其直接推断政策有效、地方服从、民众态度或单一因果。
\eventanchor{EQ-1865-DUNGAN-REVOLT-DOCUMENTS}{清廷围绕「西北回民起事在陕西、…}同治四年（1865年），清廷围绕「西北回民起事在陕西、甘肃等地扩展，形成长期战争」调遣军队、筹措饷械并处理战报，中央与地方的军政文书形成连续记录。核心取证：《清穆宗实录》同治四年相关条；军机处档、战事方略及地方奏报。这一锚点记录的是文书链：它能证明呈报、上谕、部议或照会进入机构流程，不能由此推出实际执行。来源保留：此行仅确认相关文书链和可追查的行政动作；不得由其直接推断政策有效、地方服从、民众态度或单一因果。
\eventanchor{EQ-1865-NIAN-LEADER-DEATH-DOCUMENTS}{清廷围绕「张乐行等捻军领导人先…}同治四年（1865年），清廷围绕「张乐行等捻军领导人先后失势或被杀，捻军仍持续活动」调遣军队、筹措饷械并处理战报，中央与地方的军政文书形成连续记录。核心取证：《清穆宗实录》同治四年相关条；军机处档、战事方略及地方奏报。这一锚点记录的是文书链：它能证明呈报、上谕、部议或照会进入机构流程，不能由此推出实际执行。来源保留：此行仅确认相关文书链和可追查的行政动作；不得由其直接推断政策有效、地方服从、民众态度或单一因果。
\eventanchor{EQ-1866-FUJIAN-SHIPYARD-DOCUMENTS}{围绕「福州船政局创办，成为洋务…}同治五年（1866年），围绕「福州船政局创办，成为洋务造船、教育与技术翻译的重要项目」的经费、税收、赈济或工程呈报经由户部与地方衙门处理。核心取证：《清穆宗实录》同治五年相关条；户部奏销、盐政、河工或海关档案。这一锚点记录的是文书链：它能证明呈报、上谕、部议或照会进入机构流程，不能由此推出实际执行。来源保留：此行仅确认相关文书链和可追查的行政动作；不得由其直接推断政策有效、地方服从、民众态度或单一因果。
\eventanchor{EQ-1866-NINGBO-SHANGHAI-INDUSTRY-DOCUMENTS}{围绕「江南制造局等军工项目在战…}同治五年（1866年），围绕「江南制造局等军工项目在战后江南环境中继续发展」的经费、税收、赈济或工程呈报经由户部与地方衙门处理。核心取证：《清穆宗实录》同治五年相关条；户部奏销、盐政、河工或海关档案。这一锚点记录的是文书链：它能证明呈报、上谕、部议或照会进入机构流程，不能由此推出实际执行。来源保留：此行仅确认相关文书链和可追查的行政动作；不得由其直接推断政策有效、地方服从、民众态度或单一因果。
\eventanchor{EQ-1867-NIAN-CAVALRY-CAMPAIGN-DOCUMENTS}{清廷围绕「清军以河防、堡垒和机…}同治六年（1867年），清廷围绕「清军以河防、堡垒和机动军队围堵捻军骑兵」调遣军队、筹措饷械并处理战报，中央与地方的军政文书形成连续记录。核心取证：《清穆宗实录》同治六年相关条；军机处档、战事方略及地方奏报。这一锚点记录的是文书链：它能证明呈报、上谕、部议或照会进入机构流程，不能由此推出实际执行。来源保留：此行仅确认相关文书链和可追查的行政动作；不得由其直接推断政策有效、地方服从、民众态度或单一因果。
\eventanchor{EQ-1868-NIAN-SUPPRESSION-DOCUMENTS}{清廷围绕「东捻军等主力被击败，…}同治七年（1868年），清廷围绕「东捻军等主力被击败，捻军大规模战争告终」调遣军队、筹措饷械并处理战报，中央与地方的军政文书形成连续记录。核心取证：《清穆宗实录》同治七年相关条；军机处档、战事方略及地方奏报。这一锚点记录的是文书链：它能证明呈报、上谕、部议或照会进入机构流程，不能由此推出实际执行。来源保留：此行仅确认相关文书链和可追查的行政动作；不得由其直接推断政策有效、地方服从、民众态度或单一因果。
\eventanchor{EQ-1869-YANGZHOU-SALT-CRISIS-DOCUMENTS}{围绕「两淮盐政改革、商欠与财政…}同治八年（1869年），围绕「两淮盐政改革、商欠与财政重整成为中兴时期的重要议题」的经费、税收、赈济或工程呈报经由户部与地方衙门处理。核心取证：《清穆宗实录》同治八年相关条；户部奏销、盐政、河工或海关档案。这一锚点记录的是文书链：它能证明呈报、上谕、部议或照会进入机构流程，不能由此推出实际执行。来源保留：此行仅确认相关文书链和可追查的行政动作；不得由其直接推断政策有效、地方服从、民众态度或单一因果。
\subsection{1870—1879}
\eventanchor{EQ-1870-TIANJIN-MASSACRE-DOCUMENTS}{围绕「天津教案发生，法国领事、…}同治九年（1870年），围绕「天津教案发生，法国领事、修女和中国民众死伤，引发外交危机」的照会、条约文本、换约或地方执行报告分别进入外交文书链。核心取证：《清穆宗实录》同治九年相关条；《筹办夷务始末》相关年卷与中外照会、条约文本。这一锚点记录的是文书链：它能证明呈报、上谕、部议或照会进入机构流程，不能由此推出实际执行。来源保留：此行仅确认相关文书链和可追查的行政动作；不得由其直接推断政策有效、地方服从、民众态度或单一因果。
\eventanchor{EQ-1870-ZHILI-FAMINE-RELIEF-DOCUMENTS}{围绕「华北灾荒和赈务问题促使官…}同治九年（1870年），围绕「华北灾荒和赈务问题促使官员、士绅和国际救济网络介入」的地方呈报、案件或赈务记录将社会反应带入官府文书链。核心取证：《清穆宗实录》同治九年相关条；地方志、督抚奏折及地方档案。这一锚点记录的是文书链：它能证明呈报、上谕、部议或照会进入机构流程，不能由此推出实际执行。来源保留：此行仅确认相关文书链和可追查的行政动作；不得由其直接推断政策有效、地方服从、民众态度或单一因果。
\eventanchor{EQ-1871-MUDAN-INCIDENT-DOCUMENTS}{围绕「琉球船民在台湾牡丹社一带…}同治十年（1871年），围绕「琉球船民在台湾牡丹社一带遇害，事件成为日本对台政策借口」的照会、条约文本、换约或地方执行报告分别进入外交文书链。核心取证：《清穆宗实录》同治十年相关条；《筹办夷务始末》相关年卷与中外照会、条约文本。这一锚点记录的是文书链：它能证明呈报、上谕、部议或照会进入机构流程，不能由此推出实际执行。来源保留：此行仅确认相关文书链和可追查的行政动作；不得由其直接推断政策有效、地方服从、民众态度或单一因果。
\eventanchor{EQ-1871-QING-STUDENTS-ABROAD-DOCUMENTS}{围绕「清廷派遣幼童赴美留学的计…}同治十年（1871年），围绕「清廷派遣幼童赴美留学的计划启动，洋务教育进入海外培训阶段」的禁令、编纂、教育或传播活动留下中央和地方文本记录。核心取证：《清穆宗实录》同治十年相关条；内阁档、文集或报刊相关记录。这一锚点记录的是文书链：它能证明呈报、上谕、部议或照会进入机构流程，不能由此推出实际执行。来源保留：此行仅确认相关文书链和可追查的行政动作；不得由其直接推断政策有效、地方服从、民众态度或单一因果。
\eventanchor{EQ-1872-CHINA-MERCHANTS-STEAM-NAVIGATION-DOCUMENTS}{围绕「轮船招商局创办，试图以官…}同治十一年（1872年），围绕「轮船招商局创办，试图以官督商办方式发展近代航运」的经费、税收、赈济或工程呈报经由户部与地方衙门处理。核心取证：《清穆宗实录》同治十一年相关条；户部奏销、盐政、河工或海关档案。这一锚点记录的是文书链：它能证明呈报、上谕、部议或照会进入机构流程，不能由此推出实际执行。来源保留：此行仅确认相关文书链和可追查的行政动作；不得由其直接推断政策有效、地方服从、民众态度或单一因果。
\eventanchor{EQ-1873-TONGZHI-EMPEROR-REIGN-DOCUMENTS}{围绕「同治帝亲政，垂帘体制与皇…}同治十二年（1873年），围绕「同治帝亲政，垂帘体制与皇帝个人政治角色发生调整」的上谕、奏折与部议在中央—地方行政链中流转。核心取证：《清穆宗实录》同治十二年相关条；宫中朱批奏折、内阁大库题本与《清史稿》相关志传。这一锚点记录的是文书链：它能证明呈报、上谕、部议或照会进入机构流程，不能由此推出实际执行。来源保留：此行仅确认相关文书链和可追查的行政动作；不得由其直接推断政策有效、地方服从、民众态度或单一因果。
\eventanchor{EQ-1873-JAPAN-RYUKYU-DISPUTE-DOCUMENTS}{围绕「日本围绕琉球与台湾事务加…}同治十二年（1873年），围绕「日本围绕琉球与台湾事务加强对清交涉」的照会、条约文本、换约或地方执行报告分别进入外交文书链。核心取证：《清穆宗实录》同治十二年相关条；《筹办夷务始末》相关年卷与中外照会、条约文本。这一锚点记录的是文书链：它能证明呈报、上谕、部议或照会进入机构流程，不能由此推出实际执行。来源保留：此行仅确认相关文书链和可追查的行政动作；不得由其直接推断政策有效、地方服从、民众态度或单一因果。
\eventanchor{EQ-1874-JAPANESE-TAIWAN-EXPEDITION-DOCUMENTS}{清廷围绕「日本以牡丹社事件为由…}同治十三年（1874年），清廷围绕「日本以牡丹社事件为由出兵台湾南部，清廷被迫应对」调遣军队、筹措饷械并处理战报，中央与地方的军政文书形成连续记录。核心取证：《清穆宗实录》同治十三年相关条；军机处档、战事方略及地方奏报。这一锚点记录的是文书链：它能证明呈报、上谕、部议或照会进入机构流程，不能由此推出实际执行。来源保留：此行仅确认相关文书链和可追查的行政动作；不得由其直接推断政策有效、地方服从、民众态度或单一因果。
\eventanchor{EQ-1874-TONGZHI-DEATH-DOCUMENTS}{围绕「同治帝去世，慈禧选择光绪…}同治十三年（1874年），围绕「同治帝去世，慈禧选择光绪帝继位，再次形成幼主政治」的上谕、奏折与部议在中央—地方行政链中流转。核心取证：《清穆宗实录》同治十三年相关条；宫中朱批奏折、内阁大库题本与《清史稿》相关志传。这一锚点记录的是文书链：它能证明呈报、上谕、部议或照会进入机构流程，不能由此推出实际执行。来源保留：此行仅确认相关文书链和可追查的行政动作；不得由其直接推断政策有效、地方服从、民众态度或单一因果。
